% Ho tro pdfLaTeX, XeLaTeX va LuaLaTeX tren Overleaf.
% XeLaTeX: dung Arial neu co, neu khong dung TeX Gyre Heros.
% pdfLaTeX: dung TeX Gyre Heros, font sans serif gan Helvetica/Arial.
% Khong can beamerthemeHUST.sty. Anh nen HUST trong images/ la tuy chon.
\documentclass[t,aspectratio=169,11pt]{beamer}
\usepackage{iftex}
\ifPDFTeX
  \usepackage[utf8]{inputenc}
  \usepackage[T5]{fontenc}
  \usepackage{tgheros}
  \renewcommand{\familydefault}{\sfdefault}
\else
  \usepackage{fontspec}
  \IfFontExistsTF{Arial}{
    \setmainfont{Arial}
    \setsansfont{Arial}
  }{
    \setmainfont{TeX Gyre Heros}
    \setsansfont{TeX Gyre Heros}
  }
\fi
\usefonttheme{professionalfonts}
\usepackage{booktabs,array,graphicx,tikz}
\definecolor{hustprimary}{HTML}{003563}
\definecolor{hustred}{HTML}{CE1628}
\setbeamersize{text margin left=0.65cm,text margin right=0.65cm}
\setbeamercolor{normal text}{fg=black,bg=white}
\setbeamercolor{structure}{fg=hustprimary}
\setbeamercolor{frametitle}{fg=white,bg=hustprimary}
\setbeamercolor{title}{fg=hustprimary}
\setbeamerfont{frametitle}{size=\large,series=\bfseries}
\setbeamertemplate{navigation symbols}{}
\setbeamertemplate{itemize item}{--}
\setbeamertemplate{itemize subitem}{--}
\setlength{\parskip}{0.2em}
\hypersetup{unicode=true,colorlinks=true,urlcolor=hustprimary}

% Nen goc khi co anh. Chi tai file .tex van bien dich duoc.
\newif\ifHustBackground
\newcommand{\HustBackground}[1]{%
  \HustBackgroundfalse
  \setbeamertemplate{background canvas}{}%
  \IfFileExists{images/hust-blue-169-#1.png}{%
    \HustBackgroundtrue
    \setbeamertemplate{background canvas}{%
      \includegraphics[width=\paperwidth,height=\paperheight]{images/hust-blue-169-#1.png}}%
  }{%
    \IfFileExists{hust-blue-169-#1.png}{%
      \HustBackgroundtrue
      \setbeamertemplate{background canvas}{%
        \includegraphics[width=\paperwidth,height=\paperheight]{hust-blue-169-#1.png}}%
    }{}%
  }%
}
\setbeamertemplate{frametitle}{%
  \nointerlineskip
  \begin{beamercolorbox}[wd=\paperwidth,ht=0.62cm,dp=0.2cm,leftskip=0.5cm]{frametitle}
    \usebeamerfont{frametitle}\insertframetitle
  \end{beamercolorbox}
  {\color{hustred}\hrule height 0.5pt}
  \vspace{0.25cm}
}
\setbeamertemplate{footline}{%
  \ifHustBackground\else
    \hspace*{0.65cm}{\color{hustred}\bfseries\scriptsize HUST}
  \fi
  \hfill{\color{hustprimary}\scriptsize\insertframenumber/\inserttotalframenumber}
  \hspace*{0.4cm}\vspace*{0.75cm}
}
\newcommand{\source}[1]{\par\vspace{0.25em}{\scriptsize\color{black!65}#1\par}}

\title{BÁO CÁO 2}
\subtitle{Xây dựng hệ thống thu thập, chuẩn hóa và phân tích lợi nhuận,
  dòng tiền kinh doanh trong mối liên hệ với cổ tức tiền mặt
  của doanh nghiệp niêm yết Việt Nam}
\author{Phạm Hải Nam -- 20235791\\[0.4em]
  \textbf{GVHD: TS. Nguyễn Đức Toàn}}
\date{}

\begin{document}

% 1. Bia
{\HustBackground{title}
\begin{frame}[plain]
  \begin{tikzpicture}[remember picture,overlay]
    \node[anchor=north west,xshift=0.7cm,yshift=-2.3cm,text width=11.1cm]
      at (current page.north west) {%
        {\raggedright\LARGE\bfseries\color{hustprimary}\inserttitle\par}
        \vspace{0.35cm}
        {\raggedright\large\bfseries\insertsubtitle\par}
        \vspace{0.5cm}
        {\raggedright\normalsize\insertauthor\par}
      };
  \end{tikzpicture}
\end{frame}}

\HustBackground{content-hust-full}

% 2. Boi canh
\begin{frame}{1. Bối cảnh và cơ sở đề tài}
  \begin{itemize}
    \item BCTC và thông báo cổ tức phân tán, cần ghép đúng kỳ, đơn vị,
      phạm vi báo cáo và giữ nguồn kiểm chứng.
    \item \textbf{Lợi nhuận} cần đọc cùng lợi nhuận giữ lại và các điều kiện
      phân phối, gồm nghĩa vụ tài chính, trích quỹ và thanh toán nợ.
    \item \textbf{CFO} cần đọc cùng nhu cầu đầu tư, trả nợ và tiền đầu kỳ.
    \item \textbf{Cổ tức tiền mặt} bổ sung bối cảnh về chính sách phân phối.
      Cần phân biệt năm lợi nhuận, ngày công bố và lịch trả.
  \end{itemize}
  \source{Cơ sở: Khoản 2 Điều 135 Luật Doanh nghiệp.
    CFA Institute (2026), \textit{Free Cash Flow Valuation}.
    Brav và cộng sự (2005), \textit{Payout policy in the 21st century}, mục 5.4.}
\end{frame}

% 3. Muc tieu
\begin{frame}{2. Mục tiêu và phạm vi dự kiến}
  \textbf{Mục tiêu}
  \begin{itemize}
    \item Thu thập, chuẩn hóa và kiểm tra dữ liệu từ nguồn công bố chính thức.
    \item Phân tích lợi nhuận sau thuế (LNST), dòng tiền thuần từ hoạt động
      kinh doanh (CFO) và cổ tức tiền mặt.
    \item Mở được nguồn của từng số và lưu dữ liệu để tính lại kết quả.
  \end{itemize}
  \vspace{0.3em}
\end{frame}

% 4. Chuc nang
\begin{frame}{3. Quy trình và chức năng chính}
  \begin{enumerate}
    \item \textbf{Quản lý nguồn và thu thập:} chọn doanh nghiệp, tìm và tải
      BCTC cùng thông báo cổ tức, giữ bản gốc và ngày lấy.
    \item \textbf{Trích xuất và chuẩn hóa:} tạo số ứng viên, xác định đúng
      kỳ, cột, đơn vị và phạm vi báo cáo.
    \item \textbf{Kiểm tra và duyệt:} xem số cạnh nguồn, sửa có lịch sử,
      ghi nhận phần thiếu và chặn phép tính chưa đủ dữ liệu.
    \item \textbf{Phân tích và truy nguồn:} xem biến động, cầu nối CFO,
      cổ tức theo năm lợi nhuận, mở nguồn và lưu kết quả.
  \end{enumerate}
\end{frame}

% 5. Du lieu
\begin{frame}{4. Các nhóm dữ liệu cần thu thập}
  \begin{itemize}
    \item \textbf{Doanh nghiệp và tài liệu:} mã, sàn, ngành, cơ sở quyền,
      BCTC, thông báo, bản sửa hoặc hủy, URL, kỳ, đơn vị và phiên bản.
    \item \textbf{Sáu chỉ tiêu nền:} tổng LNST, CFO, tiền và tương đương tiền,
      tổng tài sản, tổng nợ phải trả, tổng vốn chủ sở hữu.
    \item \textbf{Cầu nối và dòng tiền đã chi:} lợi nhuận trước thuế,
      các dòng điều chỉnh, CFO cuối cầu nối và tiền chi cho chủ sở hữu.
    \item \textbf{Sự kiện cổ tức:} đợt, năm lợi nhuận, tiền trên cổ phiếu,
      cơ sở quy đổi, các ngày và quan hệ sửa, hủy, trùng.
    \item \textbf{Truy nguồn và chất lượng:} trang hoặc vùng nguồn,
      trạng thái duyệt, lịch sửa và dữ liệu thiếu.
  \end{itemize}
\end{frame}

% 6. Mau va cong cu
\begin{frame}{5. Mẫu thử nghiệm VNM và công cụ}
  \textbf{BCTC hợp nhất VNM quý I/2026 đã soát xét}
  \begin{itemize}
    \item 71 trang, đơn vị VND. Cột quý I/2025 đã phân loại lại.
    \item Chọn trang PDF 12--14 để OCR, tương ứng trang in 11--13.
  \end{itemize}
  \textbf{Công cụ sử dụng}
  \begin{itemize}
    \item \texttt{urllib}, \texttt{lxml}, \texttt{hashlib}: tìm, tải và tính SHA-256.
    \item \texttt{pypdf}, Poppler: kiểm lớp chữ và chuyển trang thành ảnh.
    \item Tesseract.js và Python \texttt{re}: OCR tiếng Việt và Anh,
      trích số, đối chiếu và tính chỉ tiêu.
  \end{itemize}
  \source{\href{https://d8um25gjecm9v.cloudfront.net/cms/20260429_VNM_BCTC_DA_SOAT_XET_Q1_2026_HOP_NHAT_VN_d44326741a.pdf}{Nguồn: BCTC hợp nhất VNM quý I/2026, công bố ngày 29/04/2026.}}
\end{frame}

% 7. Ket qua cong cu
\begin{frame}{6. Kết quả thu thập và OCR ghi nhận}
  \begin{itemize}
    \item \textbf{Tìm và tải:} danh mục HTTP 200, tìm 9 liên kết PDF.
      PDF được chọn có 3.839.441 byte, lưu theo SHA-256.
    \item \textbf{Kiểm lớp chữ:} 71/71 trang trả về 0 ký tự với \texttt{pypdf}.
      Chuyển trang 12--14 thành ảnh ở 220 DPI.
    \item \textbf{OCR ba trang:} 6,026 giây
    \item \textbf{Đối chiếu sáu ô:} ghi nhận 6 đúng, 0 sai, 0 thiếu,
      so với tham chiếu đọc thủ công từ trang gốc.
  \end{itemize}

\end{frame}

% 8. So lieu va nhan xet
\begin{frame}{7. Số liệu trích xuất và nhận xét}
  {\small
  \setlength{\tabcolsep}{5pt}
  \renewcommand{\arraystretch}{1.25}
  \begin{tabular}{@{}p{3.5cm}rrp{1.1cm}@{}}
    \toprule
    \textbf{Chỉ tiêu, VND} & \textbf{Q1/2026} & \textbf{Q1/2025*} & \textbf{PDF}\\
    \midrule
    Tổng LNST hợp nhất & 2.458.221.002.532 & 1.587.273.268.054 & 12\\
    CFO & 269.326.997.516 & $-$124.187.023.735 & 13\\
    Tiền chi trả cổ tức & $-$55.837.540 & $-$1.044.977.722.500 & 14\\
    \bottomrule
  \end{tabular}}
  \vspace{0.4em}
  \begin{itemize}
    \item LNST tăng \textbf{54,87\%}. CFO tăng \textbf{393,514 tỷ đồng}, từ âm sang dương.
    \item CFO/LNST khoảng \textbf{0,110 lần}. Cần đọc tiếp các khoản điều chỉnh
      và vốn lưu động. Một quý chưa đủ kết luận chất lượng lợi nhuận.
  \end{itemize}

\end{frame}


% 10. Ket luan va buoc tiep
\begin{frame}{8. Kết luận và công việc tiếp theo}
  \small
  \textbf{Kết luận báo cáo}
  \par Đã làm rõ bài toán, mục tiêu và dữ liệu cần thu thập.
  Hệ thống hướng tới đối chiếu lợi nhuận, khả năng tạo tiền và chính sách
  cổ tức trên dữ liệu đúng kỳ, có nguồn kiểm chứng.
  \vspace{0.3em}
  \textbf{Thử nghiệm công cụ}
  \par Đã thử tìm và tải PDF VNM, kiểm lớp chữ, OCR ba trang rồi trích số.
  Sáu ô chọn trước khớp tham chiếu và tính được chỉ tiêu sơ bộ.
  Kết quả hỗ trợ tính khả thi bước đầu, chưa đại diện độ chính xác chung.
  \vspace{0.3em}
  \textbf{Công việc tiếp theo}
  \begin{itemize}
    \item Hoàn thiện trích xuất, chuẩn hóa và duyệt số có nguồn.
      Thử thêm BCTC, đo số đúng/sai/thiếu và thời gian xử lý.
    \item Tích hợp cầu nối CFO và ghép thông báo cổ tức đúng năm lợi nhuận,
      có kiểm tra bản sửa/hủy.
    \item Xây dựng giao diện đối chiếu và mở nguồn từng số.
  \end{itemize}
\end{frame}

\end{document}
